\documentclass[12pt]{article}

\usepackage{geometry}   % To control page margins
\usepackage{fancyhdr}   % To customize the header and footer
\usepackage{enumitem}   % For better lists
\usepackage{setspace}   % For line spacing
\usepackage{parskip}    % For better paragraph spacing
\usepackage{graphicx}  % Inserting images

\geometry{a4paper, margin=1in}

\pagestyle{fancy}
\fancyhf{}
\fancyhead[L]{Data Warehouses}
\fancyfoot[C]{\thepage}

\title{Overview of Data Warehousing\\\large Part 4 to the End (Slides 14--24)}
\author{0xSaad / Saad Almalki}
\date{\today}

\begin{document}

\maketitle

\section{Why OLAP and OLTP Don't Mix}

OLTP and OLAP are often separated because they have \textbf{different performance requirements}, \textbf{different data modeling requirements}, and \textbf{different data sources}.

\subsection{Different Performance Requirements}

\textbf{OLTP:}
\begin{itemize}
    \item Fast response time is important (less than 1 second in the lecture's example).
    \item Data must be up-to-date and consistent at all times.
\end{itemize}

\textbf{OLAP:}
\begin{itemize}
    \item Complex queries may consume many resources.
    \item Queries can saturate CPU and disk bandwidth.
    \item A static snapshot of data is usually acceptable.
\end{itemize}

\subsection{Different Data Modeling Requirements}

\textbf{OLTP:}
\begin{itemize}
    \item Uses \textbf{normalized schemas} for consistency.
    \item Often has complex data models and many tables.
    \item Uses a limited number of standardized queries and updates.
\end{itemize}
\textbf{OLAP:}
\begin{itemize}
    \item Simplicity of the data model is important.
    \item \textbf{Denormalized schemas} are common.
    \item Fewer joins can improve query performance.
    \item Fewer tables make the schema easier to understand.
\end{itemize}

\subsection{Analysis Requires Many Data Sources}
\begin{itemize}
    \item OLTP usually focuses on one specific process (e.g., online orders).
    \item OLAP integrates data from different processes (e.g., sales, inventory, and purchasing).
    \item OLAP often uses \textbf{historical data} to find long-term patterns and changes in behavior.
    \item Different terminology and schemas make data integration challenging.
\end{itemize}
\section{Data Warehouse as a Solution}

Using OLTP and OLAP in the same database system is often impractical. The solution is to build a \textbf{Data Warehouse}.

\begin{enumerate}
    \item Copy data from various OLTP systems.
    \item Optimize data organization and system tuning for OLAP.
    \item Prevent large analysis queries from slowing daily transactions.
    \item Periodically refresh the warehouse data.
\end{enumerate}

\section{The Warehousing Approach}

In this approach, information is \textbf{integrated in advance} and stored in a warehouse for direct querying and analysis.

\subsection{Main Architecture Components}

\begin{enumerate}
    \item \textbf{Sources:} Original systems containing the data.
    \item \textbf{Extractors/Monitors:} Extract data and monitor the sources.
    \item \textbf{Integration System:} Integrates data from different sources.
    \item \textbf{Metadata:} Descriptive information used by the integration system.
    \item \textbf{Data Warehouse:} Stores integrated information.
    \item \textbf{Clients:} Query and analyze the information.
\end{enumerate}

\textbf{Basic flow:} Sources $\rightarrow$ Extractors/Monitors $\rightarrow$ Integration System $\rightarrow$ Data Warehouse $\rightarrow$ Clients.

\subsection{Advantages}

\begin{enumerate}
    \item \textbf{High query performance}, but the information is not necessarily the most current.
    \item Complex queries at the warehouse do not interfere with local OLTP processing.
    \item Copied information can be modified, annotated, summarized, restructured, and stored historically.
\end{enumerate}

\section{Need for Data Warehousing}

Companies have gathered huge volumes of data over the years. Data warehousing helps them analyze this data and use it to gain a \textbf{competitive advantage}.

\section{Benefits of Data Warehousing}

\begin{enumerate}
    \item More efficient analysis of data.
    \item Competitive advantage.
    \item Support for strategic decision making.
    \item Increased productivity of decision makers.
\end{enumerate}

\section{Decision Support Systems (DSS), DW and OLAP}

\textbf{DSS (Decision Support System):} Information technology that helps executives, managers, and analysts make faster and better decisions.

\begin{itemize}
    \item DSS provides current and historical information for decision support.
    \item The lecture describes a \textbf{Data Warehouse} as a Decision Support System.
    \item \textbf{OLAP} is an element of DSS used for analytical processing.
\end{itemize}
\section{Data Warehouse Characteristics}

\begin{enumerate}
    \item Analysis-driven.
    \item Supports ad-hoc queries.
    \item Supports complex queries.
    \item Used by top managers.
    \item Based on dimensional modeling.
    \item Uses denormalized structures.
\end{enumerate}
\section{Data Warehousing Applications}
\begin{itemize}
    \item \textbf{Retail:} Customer loyalty and market planning.
    \item \textbf{Financial:} Risk management and fraud detection.
    \item \textbf{Airlines:} Route profitability and yield management.
    \item \textbf{Manufacturing:} Cost reduction and logistics management.
    \item \textbf{Utilities:} Asset management and resource management.
    \item \textbf{Government:} Manpower planning and cost control.
\end{itemize}

\begin{flushright}
Created by: Saad Almalki \\
\end{flushright}

\end{document}
